%! AP Statistics — Unit 8, Lesson 8.2 — Group Activity
\documentclass[10pt]{article}
\usepackage{apstats-article}
\usepackage{apstats-boxes}

\begin{document}

\pageheader{Unit 8, Lesson 8.2}{Group Activity}
\namepartnerperiod

\vspace{0.1in}
\textbf{M\&M Color Distribution Study}

The Mars company claims the following color proportions for plain M\&Ms:
Red 13\%, Orange 20\%, Yellow 14\%, Green 16\%, Blue 24\%, Brown 13\%.

You purchase a large bag and count a random sample of $n = 200$ M\&Ms.

\vspace{0.1in}
\textbf{Directions:} Complete the setup for a chi-square goodness of fit test at your tier.

\begin{tcolorbox}[colback=white, colframe=black!40,
    title=\textbf{Tier R --- Remediate},
    fonttitle=\bfseries, arc=2mm, left=3mm, right=3mm, top=2mm, bottom=2mm]

\textbf{Hypotheses} (use the sentence frames below):

H$_0$: The true color proportions for plain M\&Ms \blank{7cm}

\vspace{0.05in}

H$_a$: At least \blank{7cm}

\vspace{0.1in}
\textbf{Expected Counts.} Complete the table using $E = n \cdot p_0$ with $n = 200$.

\begin{center}
\renewcommand{\arraystretch}{1.4}
\begin{tabular}{lccc}
    \textbf{Color} & \textbf{Null proportion ($p_0$)} & \textbf{Expected count ($E = 200 \cdot p_0$)} \\
    \hline
    Red    & 0.13 & \blank{3cm} \\
    Orange & 0.20 & \blank{3cm} \\
    Yellow & 0.14 & \blank{3cm} \\
    Green  & 0.16 & \blank{3cm} \\
    Blue   & 0.24 & \blank{3cm} \\
    Brown  & 0.13 & \blank{3cm} \\
\end{tabular}
\end{center}

Are all expected counts greater than 5? \blank{2cm}

\end{tcolorbox}

\begin{tcolorbox}[colback=white, colframe=black!40,
    title=\textbf{Tier A --- Approaching Proficiency},
    fonttitle=\bfseries, arc=2mm, left=3mm, right=3mm, top=2mm, bottom=2mm]

\begin{enumerate}[label=\alph*.]

  \item State H$_0$ and H$_a$ for testing whether the M\&M colors follow the claimed
        distribution. Be sure to list all six proportions in H$_0$.

        H$_0$: \writeline

        H$_a$: \writeline

  \item Calculate the expected count for each color. Complete the table with $n = 200$.

        \begin{center}
        \renewcommand{\arraystretch}{1.4}
        \begin{tabular}{lccc}
            \textbf{Color} & \textbf{Null proportion ($p_0$)} & \textbf{Expected count ($E$)} \\
            \hline
            Red    & 0.13 & \blank{3cm} \\
            Orange & 0.20 & \blank{3cm} \\
            Yellow & 0.14 & \blank{3cm} \\
            Green  & 0.16 & \blank{3cm} \\
            Blue   & 0.24 & \blank{3cm} \\
            Brown  & 0.13 & \blank{3cm} \\
            \hline
            \textbf{Total} & 1.00 & \blank{3cm} \\
        \end{tabular}
        \end{center}

  \item Verify the Large Counts condition. Are all expected counts greater than 5?
        Show your reasoning.
        \writelines{2}

  \item What additional information would you need to check the independence condition?
        \writelines{2}

\end{enumerate}
\end{tcolorbox}

\begin{tcolorbox}[colback=white, colframe=black!40,
    title=\textbf{Tier E --- Extension},
    fonttitle=\bfseries, arc=2mm, left=3mm, right=3mm, top=2mm, bottom=2mm]

Complete all of Tier A above. Then answer the following:

\textbf{Bonus.} Design your own GOF study with a different categorical variable.

\begin{enumerate}[label=\alph*.]
  \item Describe the categorical variable and population of interest. \writelines{2}

  \item State claimed proportions for each category. Do they sum to 1? \writelines{2}

  \item Determine the minimum sample size needed so that all expected counts are at
        least 5. Show your calculation. \writelines{2}

  \item Explain why the GOF test is the appropriate procedure for your study.
        \writelines{2}
\end{enumerate}
\end{tcolorbox}

\end{document}
